\documentclass[aspectratio=169,11pt]{beamer}

\usepackage[utf8]{inputenc}
\usepackage[T1]{fontenc}
\usepackage[spanish,es-nodecimaldot]{babel}
\usepackage{booktabs}
\usepackage{array}
\usepackage{tikz}
\usepackage{amsmath}

\usetheme{Madrid}
\usecolortheme{seahorse}
\setbeamertemplate{navigation symbols}{}
\setbeamertemplate{headline}{}
\setbeamertemplate{footline}[frame number]

\definecolor{uamTeal}{RGB}{15,118,110}
\definecolor{uamOrange}{RGB}{162,62,25}
\definecolor{softGreen}{RGB}{237,247,245}
\definecolor{softOrange}{RGB}{255,244,230}

\setbeamercolor{structure}{fg=uamTeal}
\setbeamercolor{title}{fg=uamTeal}
\setbeamercolor{frametitle}{fg=uamTeal}
\setbeamercolor{block title}{bg=uamTeal,fg=white}
\setbeamercolor{block body}{bg=softGreen,fg=black}
\setbeamercolor{alerted text}{fg=uamOrange}

\newcommand{\idea}[1]{\begin{block}{Idea para recordar en 30 segundos}#1\end{block}}
\newcommand{\decision}[1]{\begin{alertblock}{Decision de ingenieria}#1\end{alertblock}}

\title{Diseno y Analisis de Experimentos en Ingenieria}
\subtitle{Presentacion didactica integrada desde practicas y casos}
\author{Roman A. Mora}
\institute{Universidad Autonoma Metropolitana}
\date{30 de junio de 2026}

\begin{document}

\begin{frame}
  \titlepage
\end{frame}

\begin{frame}{Ruta de aprendizaje}
  \idea{El curso no inicia con formulas: inicia con una pregunta sobre como decidir sin enganarse por la variabilidad.}
  \begin{enumerate}
    \item Observar un fenomeno.
    \item Convertirlo en pregunta experimental.
    \item Definir factores, niveles y respuesta.
    \item Medir con replicas y controles.
    \item Decidir con evidencia.
  \end{enumerate}
\end{frame}

\begin{frame}{Por que no basta mirar un caso?}
  \idea{Una observacion puede inspirar una pregunta, pero no demuestra una causa.}
  \begin{columns}[T]
    \begin{column}{0.48\textwidth}
      \textbf{Observacion}\\
      Aparecen hongos en tortillas bajo ciertas condiciones.
    \end{column}
    \begin{column}{0.48\textwidth}
      \textbf{Experimento}\\
      Se controlan factores, niveles, replicas y variable respuesta.
    \end{column}
  \end{columns}
  \decision{Antes de calcular, se decide que puede afirmarse y que solo puede sospecharse.}
\end{frame}

\begin{frame}{Del caso cotidiano al diseno}
  \idea{Disenar es traducir una historia en una estructura medible.}
  \begin{table}
    \centering
    \begin{tabular}{ll}
      \toprule
      Pregunta & Elemento experimental \\
      \midrule
      Que cambio? & Variable respuesta \\
      Que manipulamos? & Factores \\
      Que valores probamos? & Niveles \\
      Donde medimos? & Unidad experimental \\
      Con que comparamos? & Control o blanco \\
      \bottomrule
    \end{tabular}
  \end{table}
\end{frame}

\begin{frame}{Muestreo: el ejemplo de los gatos}
  \idea{Una muestra pequena puede contar una historia muy distinta de una muestra mas amplia.}
  \begin{center}
    \begin{tikzpicture}[scale=0.85]
      \draw[->,thick] (0,0) -- (9,0) node[right]{tamano de muestra};
      \draw[->,thick] (0,0) -- (0,3) node[above]{estabilidad};
      \draw[uamTeal,very thick] plot[smooth] coordinates{(0.4,0.6)(1.5,2.2)(2.8,1.3)(4.2,2.1)(6,2.4)(8.4,2.55)};
      \node[align=center] at (4.8,-0.65) {al aumentar $n$, la percepcion se vuelve menos caprichosa};
    \end{tikzpicture}
  \end{center}
\end{frame}

\begin{frame}{ANOVA como comparacion de variabilidades}
  \idea{ANOVA pregunta si la variacion entre tratamientos es grande frente a la variacion natural dentro de ellos.}
  \[
    F=\frac{\text{variabilidad entre tratamientos}}{\text{variabilidad dentro de tratamientos}}
  \]
  \decision{Un valor p no sustituye la interpretacion: solo ayuda a decidir si la diferencia observada merece confianza estadistica.}
\end{frame}

\begin{frame}{Lentejas: miel, azucar, cafe y blanco}
  \idea{La practica de lentejas permite aprender ANOVA sin perder de vista el sistema biologico.}
  \begin{table}
    \centering
    \begin{tabular}{lll}
      \toprule
      Factor & Niveles & Respuestas \\
      \midrule
      Aditivo & miel, azucar, cafe, blanco & germinacion, raiz, altura, hongos \\
      Tiempo & dias de observacion & cambio del sistema \\
      \bottomrule
    \end{tabular}
  \end{table}
\end{frame}

\begin{frame}{Dos factores: no todo se suma}
  \idea{Una interaccion aparece cuando el efecto de un factor depende del nivel del otro.}
  \begin{columns}
    \begin{column}{0.5\textwidth}
      \textbf{Efecto principal}\\
      Un factor cambia la respuesta en promedio.
    \end{column}
    \begin{column}{0.5\textwidth}
      \textbf{Interaccion}\\
      La combinacion produce algo que no se explica por separado.
    \end{column}
  \end{columns}
  \decision{Si hay interaccion, interpretar solo promedios generales puede ocultar el fenomeno importante.}
\end{frame}

\begin{frame}{Factorial $2^2$: curcuma y cloro en papel}
  \idea{Cuatro tratamientos bien definidos pueden revelar efectos principales e interacciones.}
  \begin{table}
    \centering
    \begin{tabular}{ccc}
      \toprule
      Tratamiento & Curcuma & Cloro \\
      \midrule
      1 & bajo & bajo \\
      2 & alto & bajo \\
      3 & bajo & alto \\
      4 & alto & alto \\
      \bottomrule
    \end{tabular}
  \end{table}
  \decision{El blanco y las replicas protegen contra conclusiones vistosas pero fragiles.}
\end{frame}

\begin{frame}{Cuando la realidad no queda balanceada}
  \idea{En procesos manuales o industriales, el rigor consiste en documentar restricciones y no esconderlas.}
  \begin{itemize}
    \item replicas desiguales;
    \item variabilidad del operador;
    \item materiales no identicos;
    \item mediciones con error.
  \end{itemize}
  \decision{Un diseno imperfecto puede ser util si sus limites quedan claros.}
\end{frame}

\begin{frame}{Taguchi: explorar sin probarlo todo}
  \idea{Una matriz reducida ayuda cuando el sistema tiene varios factores y recursos limitados.}
  \begin{table}
    \centering
    \begin{tabular}{cccc}
      \toprule
      Corrida & Cafe & Miel & Azucar \\
      \midrule
      1 & bajo & bajo & bajo \\
      2 & bajo & alto & alto \\
      3 & alto & bajo & alto \\
      4 & alto & alto & bajo \\
      \bottomrule
    \end{tabular}
  \end{table}
\end{frame}

\begin{frame}{Violeta africana: cuando el resultado ensena otra cosa}
  \idea{Si no aparece la propagacion esperada, el experimento no fracasa automaticamente: puede revelar nuevas respuestas.}
  \begin{itemize}
    \item cobertura fungica;
    \item halo cafe;
    \item tejido verde;
    \item necrosis.
  \end{itemize}
  \decision{El analisis debe distinguir hallazgo exploratorio de conclusion definitiva.}
\end{frame}

\begin{frame}{ANCOVA y MANOVA}
  \idea{No todos los problemas se resuelven comparando una sola media.}
  \begin{columns}[T]
    \begin{column}{0.48\textwidth}
      \textbf{ANCOVA}\\
      Compara tratamientos ajustando por covariables.
    \end{column}
    \begin{column}{0.48\textwidth}
      \textbf{MANOVA}\\
      Evalua varias respuestas relacionadas como un sistema.
    \end{column}
  \end{columns}
\end{frame}

\begin{frame}{Proyecto integrador}
  \idea{El cierre del curso es defender una decision experimental completa.}
  \begin{enumerate}
    \item Problema de ingenieria.
    \item Pregunta e hipotesis.
    \item Factores, niveles y respuesta.
    \item Procedimiento y datos.
    \item Analisis e interpretacion.
    \item Limitaciones y siguiente experimento.
  \end{enumerate}
\end{frame}

\begin{frame}{Criterio final}
  \idea{Un buen experimento no promete certeza absoluta: produce evidencia organizada para decidir mejor.}
  \decision{La pregunta final no es solo ``salio significativo'', sino: que puedo afirmar, con que limite y cual seria el siguiente experimento?}
\end{frame}

\end{document}
